\chapter{CONSIDERAÇÕES FINAIS}

Este projeto buscou estruturar os caminhos para uma pesquisa que entende o espaço geográfico como uma instância social e historicamente construída. Ao longo das etapas previstas, espera-se demonstrar que o ensino de Geografia, ancorado nas reflexões de autores como \cite{milton-ne}, pode ser uma ferramenta poderosa para a formação de sujeitos críticos.

As considerações finais serão tecidas ao término da execução da pesquisa, mas já apontam para a importância de práticas educativas que dialoguem com a realidade socioespacial dos alunos e superem os modelos tradicionais de avaliação, conforme discutido por \cite{luckesi-avaprend}.